The backend milestone was where everything became connected.
The frontend visualisation linked to the backend API endpoint, labelled /getGraph (alongside the necessary request parameters),
to show data that was actually relevant to the user,
rather than placeholder mock data.
Due to the gravity of this milestone, it is, of course, split into several stages:
the power of Spring Boot, implementing the API endpoint itself, and (unit) testing.


\section{The power of Spring Boot}\label{sec:the-power-of-spring-boot}

Without Spring Boot, there would have been a lot more work in implementing the API endpoint.
Spring Boot was incredibly powerful in terms of what was achieved with such little code.
All that was needed was the instantiation of two list objects
(two ArrayList objects, in particular) with particular attributes,
and Spring Boot handled the rest of the conversion to a JSON object.
Of course,
adding specific values to these list objects was more challenging and complex as explained in the Frontend chapter.
A single controller class controlled the entire endpoint,
and there was no concern about this
since there was no sensitive data being transferred other than publicly available information and only a single endpoint was being developed for the system.
Hence, request parameters, within the URL of the endpoint itself, were used to pass the owner of the project and the project name.
A start and end timestamp were optional parameters, implemented through two Material UI Date Time Pickers in the frontend system.
From the use of Java functions to perform endpoint actions, this made testing very straightforward later on.


\section{Implementing the API endpoint itself}\label{sec:implementing-the-api-endpoint-itself}

Implementing the actual API endpoint itself was easier said than done, primarily due to the complex nature of the algorithm powering the relationship between developers.
For the algorithm to be implemented properly, three new objects had to be created in the backend: Node, Link, and Graph.
The Node class stored author information for every pull request that was stored in the pull requests table in the relational database alongside a numerical size value.
This numerical size value was determined based on the number of comments an author has made on the whole project.
The Link class stored the associated developers who communicated through a comment(s) on the SAME pull request, containing a source and target node alongside a numerical thickness value.
This numerical thickness value is based on the number of comments made between two developers on the SAME pull request(s).
Finally, the Graph class stored a list of Nodes and Links respectively in the form of two ArrayList objects.
This Graph class is what was actually returned by the API endpoint, resulting in the carefully structured JSON response.


\section{Testing}\label{sec:testing}

Testing was of lower priority when developing the backend system of the project, however, it was still an important aspect.
Regardless of whether users only viewed the visualisation, it was still paramount that all tests passed with adequate test coverage.
It was discovered early on in the testing process that it would be inefficient and unwise to perform tests based on live production database information.
This was because the database information could be removed/changed at any time.
Subsequently, an in-memory database solution, H2 in particular, was applied for this purpose.
The primary reason why was chosen was for that very reason: H2 is in-memory only.
This means that the contents of the database are wiped between test runs, and thus helps maintain a consistent testing environment.
It also means that the contents of the tests can be manually chosen and not reliant on live production data.
In the end, all tests passed successfully, with twenty eight tests being run in total, with 97\% overall test coverage.